\documentclass[11pt]{article}

\usepackage[margin=1in]{geometry}
\usepackage{amsmath,amssymb,mathtools}
\usepackage{booktabs}
\usepackage{array}
\usepackage{hyperref}

\title{HAOS-IIP Operator Architecture and Low-Spectrum Results}
\author{HAOS-IIP Research Stack}
\date{March 10, 2026}

\begin{document}
\maketitle

\begin{abstract}
This note summarizes the current HAOS-IIP operator program at the level that is now supported by explicit numerics. A single interaction kernel defines a weighted discrete substrate, but different operator lifts act on different field types. The node Laplacian $L_0$ produces the scalar or geometry branch. The edge or Hodge operator $L_1$ produces the first non-scalar branch. The present HAOS/IIP edge operator supports a genuine non-scalar flux-responsive branch, but explicit harmonic-vs-coexact separation shows that the lowest modes remain harmonic or mixed topological structure rather than a low coexact vector band. In the current periodic/twisted range, this is a negative Maxwell test. No claim of emergent photons, full gauge theory, or particle matching is made here.
\end{abstract}

\section{Introduction}
The current HAOS-IIP architecture is most cleanly stated as
\begin{equation}
\text{same substrate} \quad + \quad \text{same kernel} \quad + \quad \text{different operator lifts}.
\end{equation}
The project no longer treats a single node Laplacian as the entire physical content. Instead, the present operator stack has three distinct levels:
\begin{align}
L_0 &: \text{node operator for scalar and geometry-like structure}, \\
L_1 &: \text{edge operator for vector-sector tests}, \\
D_H &: \text{Dirac-type lift for later fermion-sector work}.
\end{align}
Only the first two levels are numerically active in the present repository. Operationally, interaction-invariant physics means that these different operator lifts are built on the same kernel-defined weighted substrate rather than on unrelated microscopic models. The node sector is already supported by explicit low-mode studies on concrete $3$D substrates. The edge sector now has a working periodic and twisted implementation together with a wider flux scan. The main question addressed in this paper is therefore narrow: does the current edge operator contain a genuine low non-scalar branch, or is the apparent vector content still inherited from scalar gradients?

\begin{figure}[t]
\centering
\small
\setlength{\fboxsep}{6pt}
\renewcommand{\arraystretch}{1.1}
\begin{tabular}{c c c c c c c c c}
\fbox{\parbox{0.105\linewidth}{\centering interaction\\kernel}} &
$\to$ &
\fbox{\parbox{0.105\linewidth}{\centering weighted\\graph}} &
$\to$ &
\fbox{\parbox{0.105\linewidth}{\centering $L_0$\\node sector}} &
$\to$ &
\fbox{\parbox{0.105\linewidth}{\centering $L_1$\\edge sector}} &
$\to$ &
\fbox{\parbox{0.105\linewidth}{\centering $D_H$\\Dirac lift}}
\end{tabular}
\caption{Current HAOS-IIP operator hierarchy. One interaction kernel defines the weighted substrate, while different operator lifts target scalar, vector, and later spinorial sectors. Only $L_0$ and $L_1$ are numerically active in the present repository.}
\label{fig:operator-hierarchy}
\end{figure}

\section{Interaction Kernel}
The common starting point is a weighted interaction kernel on embedded nodes $x_i \in \mathbb{R}^3$,
\begin{equation}
A_{ij} = \exp\!\left(-\frac{|x_i-x_j|^2}{2\varepsilon_k}\right),
\qquad
A_{ii}=0.
\end{equation}
Here $\varepsilon_k$ is the kernel width. For numerics, the graph may be truncated at a finite interaction radius
\begin{equation}
r_c = \alpha \sqrt{\varepsilon_k},
\end{equation}
with the truncation factor $\alpha$ reported explicitly.

The weighted graph data are
\begin{equation}
D_{ii} = \sum_j A_{ij},
\qquad
L_0 = D - A.
\end{equation}
For the edge complex, the same kernel determines edge weights. On a nearest-neighbor cubic complex with edge length $h$,
\begin{equation}
w_e = \exp\!\left(-\frac{h^2}{2\varepsilon_k}\right).
\end{equation}
The key point is that the scalar and edge sectors are not built from unrelated models. They are different lifts of the same weighted interaction fabric.

\section{Graph Substrate}
Two substrate classes are currently used.

First, for the scalar study, nodes are embedded either as a perturbed cubic lattice or as a random geometric point cloud in $[0,1]^3$. Edges are included when $|x_i-x_j| \le r_c$, with weight $A_{ij}$.

Second, for the edge study, the substrate is a weighted cubic cell complex. Let $V$, $E$, and $F$ denote nodes, oriented edges, and oriented faces. The node-edge incidence matrix $B_0$ is defined by
\begin{equation}
(B_0)_{e,i} =
\begin{cases}
-1, & i = \mathrm{tail}(e), \\
+1, & i = \mathrm{head}(e), \\
0, & \text{otherwise},
\end{cases}
\end{equation}
and the face-edge incidence matrix $C$ records oriented face boundaries.

For the open cubic complex used in the first $L_1$ tests, the weighted differentials are
\begin{equation}
d_0 = W_1^{1/2} B_0,
\qquad
d_1 = W_2^{1/2} C,
\end{equation}
with diagonal edge and face weight matrices $W_1$ and $W_2$.

For the periodic gauge-sector study, opposite faces of the cubic box are identified to form a discrete $3$-torus $\mathbb{T}^3_n$. This periodic identification is not cosmetic. It introduces a nontrivial low harmonic sector on edge space that is absent on the open box.

\section{Node Operator L0}
The scalar operator acts on node $0$-forms,
\begin{equation}
(L_0 \psi)_i = \sum_j A_{ij}(\psi_i-\psi_j).
\end{equation}
When a graph incidence representation is available, one may also write
\begin{equation}
L_0 = B_0^\top W_1 B_0.
\end{equation}
This is the branch that currently carries the geometry evidence.

The strongest scalar benchmark in the repository is the open perturbed cubic study on a $6 \times 6 \times 6$ substrate with $N=216$ nodes and reference kernel width $\varepsilon_k = 0.056997$. The first nontrivial eigenvalues are
\begin{equation}
\lambda_1 = 2.410405,
\qquad
\lambda_2 = 2.415782,
\qquad
\lambda_3 = 2.422536.
\end{equation}
These three modes form a near-degenerate smooth family aligned with the coordinate directions of the box. On a nearly cubic domain, the first nontrivial scalar Laplacian modes are expected to vary primarily along the three coordinate axes, so this near-triple degeneracy is the discrete analogue of the first box-mode family. Their inverse participation ratios remain low, and the visual mode shapes match the first scalar box modes. This is the current evidence that the node branch reproduces geometry-like long-wavelength behavior.

The same low-spectrum study also tested an open random geometric graph. There the lowest modes were more strongly mixed with localization and boundary artifacts. That comparison is important because it sets the present diagnostic choice: regular or weakly perturbed cubic substrates are the cleanest setting for operator separation, whereas random graphs mix sectors too early.

A phase-dressed scalar extension has also been implemented,
\begin{equation}
(L_\theta \psi)_i = \sum_j A_{ij}(\psi_i - U_{ij}\psi_j),
\qquad
U_{ij} = e^{i\theta_{ij}}.
\end{equation}
This branch is connection sensitive and supports organized edge currents, but it still acts on node scalars. It is therefore a bridge toward gauge structure, not the gauge sector itself.

\section{Edge Operator L1}
The vector-sector candidate acts on oriented edge fields. On the weighted cubic complex, the implemented edge operator is the discrete Hodge Laplacian
\begin{equation}
L_1 = d_0 d_0^* + d_1^* d_1.
\end{equation}
In the open real-valued implementation,
\begin{equation}
d_0 = W_1^{1/2} B_0,
\qquad
d_1 = W_2^{1/2} C,
\qquad
L_1 = d_0 d_0^\top + d_1^\top d_1.
\end{equation}
The edge space decomposes as
\begin{equation}
C^1 = \operatorname{im} d_0 \oplus \mathcal{H}^1 \oplus \operatorname{im} d_1^*,
\end{equation}
where exact components lie in $\operatorname{im} d_0$, coexact components lie in $\operatorname{im} d_1^*$, and harmonic modes lie in
\begin{equation}
\mathcal{H}^1 = \ker d_0^* \cap \ker d_1.
\end{equation}
On a contractible open box, the low harmonic sector is trivial. On a $3$-torus, it is not.

\begin{figure}[t]
\centering
\small
\setlength{\fboxsep}{6pt}
\begin{tabular}{c c c c c}
\fbox{\parbox{0.18\linewidth}{\centering exact\\$\operatorname{im} d_0$}} &
$\oplus$ &
\fbox{\parbox{0.18\linewidth}{\centering harmonic\\$\mathcal{H}^1$}} &
$\oplus$ &
\fbox{\parbox{0.18\linewidth}{\centering coexact\\$\operatorname{im} d_1^*$}}
\end{tabular}
\caption{Edge-space decomposition used to interpret the low $L_1$ spectrum. On the periodic branch, the divergence-free subspace contains both harmonic and coexact components.}
\label{fig:hodge-decomposition}
\end{figure}

The first open-box $L_1$ study showed that the bottom of the edge spectrum is exact dominated. For the open $n=5$ cubic box, the first three eigenvalues are all
\begin{equation}
0.326713,
\end{equation}
and each corresponding mode has exact fraction $1.000$ and coexact fraction $0.000$. In other words, the unrestricted low edge spectrum of the open box is still scalar leakage.

This is the baseline that the periodic and twisted experiment had to beat.

\section{Periodic Twisted Experiment}
The next step replaces the open box by a periodic cubic lattice of side length $n$ with nearest-neighbor oriented edges in the $x$, $y$, and $z$ directions. The edge weight remains kernel derived,
\begin{equation}
w_e = \exp\!\left(-\frac{h^2}{2\varepsilon_k}\right),
\qquad
h = \frac{1}{n},
\qquad
\varepsilon_k = 0.2.
\end{equation}
A background $U(1)$ connection is added in a discrete Landau gauge with flux quantum $m$ through each $xy$ plaquette,
\begin{equation}
\phi = \frac{2\pi m}{n^2}.
\end{equation}
Here $m$ counts the inserted plaquette-flux quanta through the periodic $xy$ cycles, so increasing $m$ twists the background connection and lifts the torus-cycle zero modes.
The implemented link phases are
\begin{align}
U_x(i,j,k) &=
\begin{cases}
1, & i < n-1, \\
\exp(-i\phi n j), & i = n-1,
\end{cases} \\
U_y(i,j,k) &= \exp(i\phi i), \\
U_z(i,j,k) &= 1.
\end{align}
The covariant node-edge differential is
\begin{equation}
(d_{0,A}\psi)_e = \sqrt{w_e}\bigl(U_e\psi_{\mathrm{head}(e)} - \psi_{\mathrm{tail}(e)}\bigr).
\end{equation}
A compatible covariant face operator $d_{1,A}$ is built by transporting boundary edge values back to a chosen face basepoint before summing the oriented boundary. The twisted edge operator is then
\begin{equation}
L_{1,A} = d_{0,A} d_{0,A}^* + d_{1,A}^* d_{1,A}.
\end{equation}

For the periodic analysis, exact and divergence-free content are measured using orthogonal projectors
\begin{equation}
P_{\mathrm{ex}} : C^1 \to \operatorname{im} d_{0,A},
\qquad
P_{\mathrm{div}} : C^1 \to \ker d_{0,A}^*.
\end{equation}
For a mode $a$, the reported fractions are
\begin{equation}
f_{\mathrm{ex}}(a) = \frac{\|P_{\mathrm{ex}} a\|^2}{\|a\|^2},
\qquad
f_{\mathrm{div}}(a) = \frac{\|P_{\mathrm{div}} a\|^2}{\|a\|^2}.
\end{equation}
Because $\ker d_{0,A}^*$ contains both coexact and harmonic pieces, the reported ``coexact fraction'' in the experiment should be read operationally as a divergence-free fraction.

The open and periodic branches then separate sharply.

For the open $n=5$ box, the first three $L_1$ modes remain exact gradients:
\begin{center}
\begin{tabular}{@{}ccccc@{}}
\toprule
mode & $\lambda$ & exact & divergence-free & pattern \\
\midrule
0 & 0.326713 & 1.000 & 0.000 & exact $x$-gradient \\
1 & 0.326713 & 1.000 & 0.000 & exact $y$-gradient \\
2 & 0.326713 & 1.000 & 0.000 & exact $z$-gradient \\
\bottomrule
\end{tabular}
\end{center}

For the periodic $n=5$ torus at trivial flux $m=0$, the lowest three modes are instead harmonic cycles:
\begin{center}
\begin{tabular}{@{}ccccc@{}}
\toprule
mode & $\lambda$ & exact & divergence-free & pattern \\
\midrule
0 & 0.000000 & 0.000 & 1.000 & harmonic $y$-cycle \\
1 & 0.000000 & 0.000 & 1.000 & harmonic $x$-cycle \\
2 & 0.000000 & 0.000 & 1.000 & harmonic $z$-cycle \\
\bottomrule
\end{tabular}
\end{center}

At nontrivial flux $m=1$, those zero modes lift while staying mostly divergence free:
\begin{center}
\begin{tabular}{@{}cccccc@{}}
\toprule
mode & $\lambda$ & exact & divergence-free & $\|d_{0,A}^* a\|$ & $\|d_{1,A} a\|$ \\
\midrule
0 & 0.013057 & 0.025 & 0.975 & 0.0850 & 0.0764 \\
1 & 0.111287 & 0.194 & 0.806 & 0.2344 & 0.2373 \\
2 & 0.220336 & 0.000 & 1.000 & 0.0000 & 0.4694 \\
3 & 0.221528 & 0.097 & 0.903 & 0.3099 & 0.3542 \\
\bottomrule
\end{tabular}
\end{center}

The divergence-free projected spectrum for the same case begins at
\begin{equation}
0.027359,\ 0.220336,\ 0.243475,\ 0.440293, \dots
\end{equation}
This matters because the low branch survives projection onto $\ker d_{0,A}^*$ and therefore cannot be reduced to exact scalar gradients, even though this projector does not by itself separate harmonic and coexact content.

The same qualitative behavior persists for $n=4$. For $m=1$, the first mode has
\begin{equation}
\lambda = 0.014110,
\qquad
f_{\mathrm{ex}} = 0.020,
\qquad
f_{\mathrm{div}} = 0.980.
\end{equation}
The low periodic branch is therefore not a single-size artifact.

\section{Flux Scan Results}
The wider flux scan extends the periodic and twisted experiment to
\begin{equation}
m = 0,1,2,3,4
\end{equation}
for lattice sizes $n=4$ and $n=5$. For each case, the code records low eigenvalues, exact fraction, divergence-free fraction, divergence norm, curl norm, inverse participation ratio, and qualitative support pattern.

To follow one representative low branch, the scan selects the lowest mode among the first four full-spectrum modes whose divergence-free fraction is at least $0.8$. This is an operational diagnostic, not a proof of a continuum band.

For $n=5$, the selected low branch is
\begin{center}
\begin{tabular}{@{}cccc@{}}
\toprule
$m$ & $\phi$ & $\lambda$ & divergence-free fraction \\
\midrule
0 & 0.000000 & 0.000000 & 1.0000 \\
1 & 0.251327 & 0.013057 & 0.9755 \\
2 & 0.502655 & 0.018837 & 0.9837 \\
3 & 0.753982 & 0.024179 & 0.9876 \\
4 & 1.005310 & 0.032351 & 0.9870 \\
\bottomrule
\end{tabular}
\end{center}
This branch is nearly smooth from $m=1$ onward.

For $n=4$, the selected low branch is
\begin{center}
\begin{tabular}{@{}cccc@{}}
\toprule
$m$ & $\phi$ & $\lambda$ & divergence-free fraction \\
\midrule
0 & 0.000000 & 0.000000 & 1.0000 \\
1 & 0.392699 & 0.014110 & 0.9798 \\
2 & 0.785398 & 0.171805 & 0.9228 \\
3 & 1.178097 & 0.049192 & 0.9821 \\
4 & 1.570796 & 0.275223 & 0.9077 \\
\bottomrule
\end{tabular}
\end{center}
The divergence-free content remains high, but the spectral ordering is less smooth.

The projected spectrum confirms the same pattern. The first nonzero projected eigenvalue evolves as
\begin{align}
&n=5: && 0.000000 \to 0.027359 \to 0.035496 \to 0.041094 \to 0.058335, \\
&n=4: && 0.000000 \to 0.032130 \to 0.197545 \to 0.083671 \to 0.327827.
\end{align}
The $n=5$ branch is cleaner; the $n=4$ branch still shows strong reordering. The broad conclusions of the scan are therefore:
\begin{enumerate}
\item the low periodic branch stays mostly divergence free across the scanned flux values,
\item the branch is flux responsive rather than collapsing back into the exact sector,
\item the behavior persists across the tested lattice sizes,
\item the low branch does not yet organize into one clean smooth transverse band across the full scan.
\end{enumerate}

A stricter follow-up test now separates the untwisted periodic edge space into exact, harmonic, and coexact sectors and then restricts the twisted operator to the coexact subspace. Current gauge-sector verdict: The present HAOS/IIP edge operator supports a genuine non-scalar flux-responsive branch, but explicit harmonic-vs-coexact separation shows that the lowest modes remain harmonic or mixed topological structure rather than a low coexact vector band. In the current periodic/twisted range, this is a negative Maxwell test. Quantitatively,
\begin{itemize}
\item $n=5$, $m=1$: full lowest $0.013057$, harmonic-projected $0.509808$, coexact-projected $0.458457$
\item $n=4$, $m=1$: full lowest $0.014110$, harmonic-projected $0.611897$, coexact-projected $0.679582$
\end{itemize}
So the projected coexact sector exists and is numerically stable, but it does not form the actual low band once harmonic torus cycles are removed.

\section{Interpretation}
The current numerical evidence supports the kinematics of an operator-level vector sector, but not a full gauge theory.

The decisive point is the contrast between open and periodic behavior. On the open box, the bottom of $L_1$ is exact dominated and therefore reducible to scalar leakage. On the periodic torus, a low non-exact and divergence-free branch appears and responds to nontrivial holonomy. This is explicit non-scalar structure in the edge operator.

The present implementation also has an important limitation. The earlier divergence-free projector onto $\ker d_{0,A}^*$ mixes harmonic and coexact content on the torus. The explicit Hodge-separation follow-up resolves that ambiguity. After projecting out both the exact sector and the harmonic torus cycles, the coexact sector remains real and numerically stable, but it does not become the actual low band. In the tested cases the coexact floor stays much higher than the mixed full low branch.

This places the current HAOS-IIP stack close to the kinematic side of lattice gauge theory. The repository already contains
\begin{enumerate}
\item $L_0 = d_0^* d_0$,
\item $L_1 = d_0 d_0^* + d_1^* d_1$,
\item link phases $U_{ij}$,
\item plaquette holonomy and flux response,
\item explicit low non-scalar edge modes.
\end{enumerate}
What it does not yet contain is the dynamical part of a gauge theory:
\begin{enumerate}
\item a dynamical action for the link variables $U$,
\item a plaquette energy term playing the role of $F^2$,
\item a clean propagating transverse band beyond harmonic torus cycles.
\end{enumerate}

The correct reading is therefore limited but meaningful:
\begin{quote}
The repository now contains operator-level evidence for a non-scalar vector-like branch in $L_1$. That branch is flux responsive on topologically nontrivial substrates, but explicit Hodge separation shows that the lowest modes remain harmonic or mixed topological structure rather than a low coexact vector band.
\end{quote}
This distinction matters. The numerics now justify talking about a non-scalar edge sector. They do not justify talking about emergent photons, a derived Maxwell action, or physical couplings.

\section{Open Problems}
The present operator architecture is incomplete in several concrete ways.

First, the explicit Hodge split has now been implemented, and the result is a negative Maxwell test in the current range: the coexact floor does not descend into the actual low band. The next problem is therefore not to redefine the decomposition, but to find out whether the spectrum can be reorganized so that local circulation modes descend below the harmonic or mixed topological floor.

Second, the periodic torus is not sufficient by itself. Punctures and controlled defect backgrounds are needed to break the dominance of global torus cycles and test whether more local vector excitations appear lower in the spectrum.

Third, larger lattice sizes are required to determine whether the coexact floor moves down toward a cleaner band or simply stabilizes as a higher sector.

Fourth, the relation between the phase-dressed scalar bridge $L_\theta$ and the edge branch $L_{1,A}$ still needs continuum control. At present, the code demonstrates connection sensitivity and flux response, but not a derivation of a dynamical gauge field from the kernel alone.

Fifth, the fermion branch remains postponed. A Dirac-type operator $D_H$ is part of the planned architecture, but it should only be interpreted after the gauge branch is cleaner. Otherwise any spinorial numerics will be hard to distinguish from artifacts of an incompletely resolved connection sector.

The immediate numerical targets are therefore:
\begin{enumerate}
\item puncture and defect backgrounds,
\item larger periodic sizes,
\item continuum scaling tests for the coexact floor,
\item only then a Dirac-type lift on the same weighted substrate or a dynamical plaquette-action branch.
\end{enumerate}

\section{Conclusion}
The present HAOS-IIP results support a clean operator architecture built from one weighted interaction substrate and multiple operator lifts. The node operator $L_0$ already reproduces a scalar or geometry-like low spectrum on regular $3$D substrates. The edge operator $L_1$ now passes the main falsification test that had remained open: on a periodic and twisted cubic substrate, the low branch is not reducible to exact scalar gradients. It remains mostly divergence free, responds to flux, and persists across the tested sizes.

That is the current positive result. The current limitation is equally clear: the same low branch is still best interpreted as flux-lifted harmonic or topological edge structure rather than a clean propagating Maxwell-like band. The project therefore has explicit operator-level evidence for a vector sector, but not yet a full emergent gauge theory.

\end{document}
